\PassOptionsToPackage{table,xcdraw}{xcolor}
\documentclass[10pt, a4paper, landscape]{article}

\usepackage[T1]{fontenc}
\usepackage[utf8]{inputenc}
\usepackage{tgpagella}
\usepackage[spanish, es-noshorthands]{babel}
\usepackage{amsmath, amssymb}
\usepackage{booktabs}
\usepackage{tabularx}
\usepackage[table]{xcolor}
\usepackage[margin=1.4cm]{geometry}
\usepackage{enumitem}
\usepackage{microtype}
\usepackage{hyperref}

\emergencystretch=2em
\tolerance=1000
\hbadness=10000

\definecolor{ucbblue}{RGB}{0, 51, 102}
\definecolor{ucbgray}{RGB}{240, 244, 248}
\definecolor{warnred}{RGB}{180, 20, 20}

\begin{document}

\noindent
{\LARGE \textbf{\textcolor{ucbblue}{IMT-221 — Circuitos Electrónicos III}}}\\[0.12cm]
{\Large \textbf{Hito 3: Guía Única E1 $\rightarrow$ E4 — Parte 5: Plan de Mediciones, Plantilla, Puertas G1--G4 y Entrega Final (v1.5)}}\\[0.08cm]
{\small \textbf{Versión integrada 1.5} $\cdot$ 12 de noviembre de 2026 $\cdot$ Protocolo de Auditoría, Entrega Incremental y Fuentes}
\par\vspace{0.15cm}\noindent\rule{\linewidth}{0.4pt}\vspace{0.2cm}

\section*{7. Plan de Mediciones Unificado (Aplicable a Todas las Etapas)}

Completar una tabla por etapa y condición, sin borrar nunca una fila pendiente. El informe de cada equipo contendrá al menos:

\par\vspace{0.1cm}\noindent
\renewcommand{\arraystretch}{1.25}
\begin{tabularx}{\linewidth}{|c|>{\raggedright\arraybackslash}p{6.4cm}|>{\raggedright\arraybackslash}p{6.2cm}|>{\raggedright\arraybackslash}X|}
\hline
\rowcolor{ucbgray}
\textbf{Etapa} & \textbf{Señales o Datos que se Guardan} & \textbf{Variable Controlada} & \textbf{Comparación Obligatoria} \\ \hline
\textbf{E1} & $V_{\text{BUS}}$, $V_{GS}$, $V_{DS}$, $I_D$, $f$, duty, $R_L$ y calor & $0\text{--}5\,\text{V}$ vs.\ $0\text{--}10\,\text{V}$ según MOSFET y fuente & Ley de Ohm, conducción, nivel de gate \\ \hline
\textbf{E2} & \texttt{PWM\_A}, \texttt{OPTO\_N}, \texttt{VIN\_ISO}, \texttt{BH}, \texttt{LCTL}, $V_{GS}$, $V_{DS}$, $I_D$, $f$, duty & E1 directo frente a E2 optoaislado con mismo MOSFET/carga & Nivel alto/bajo, carga/descarga y pérdidas \\ \hline
\textbf{E3} & $V_{GS}$, $V_{DS}$, $i_L$, $i_{DF}$, picos, temperatura & Sin/con protección \textbf{en simulación}; protección real aprobada & Sobrevoltajes, continuidad de corriente, energía \\ \hline
\textbf{E4} & $V_{GS}$ de 4 MOSFET, nodos A/B, $I_{\text{MOTOR}}$, $I_{\text{BUS}}$, $V_{\text{BUS}}$, tiempos muertos & Secuencia y frecuencia PWM aprobadas & Estados, no solapamiento, pérdidas/térmica \\ \hline
\end{tabularx}

\subsection*{Plantilla de Registro Trazable (Repetir por Magnitud)}
\par\vspace{0.08cm}\noindent
\renewcommand{\arraystretch}{1.28}
{\small
\begin{tabularx}{\linewidth}{|>{\raggedright\arraybackslash}p{7.8cm}|>{\raggedright\arraybackslash}X|}
\hline
\rowcolor{ucbgray}
\textbf{Campo} & \textbf{Anotación} \\ \hline
Equipo / integrantes / integrante que explica & \\ \hline
Fecha / etapa / G1--G4 & \\ \hline
Pieza real, fabricante, datasheet y encapsulado & \\ \hline
Instrumento, modelo, sonda/escala, conexión de referencia & \\ \hline
Condición de tensión, carga, frecuencia y duty & \\ \hline
Magnitud y unidad & \\ \hline
\texttt{CALC} — Valor y fórmula & \\ \hline
\texttt{SIM} — Valor, modelo, archivo, log y captura & \\ \hline
\texttt{MED} — Valor y captura (\textbf{solo si medido}) & \\ \hline
Diferencia absoluta/relativa si es significativa & \\ \hline
Incertidumbre/resolución y posible explicación & \\ \hline
Pendiente, anomalía o condición de parada & \\ \hline
\end{tabularx}
}

\newpage
\section*{8. Puertas de Seguridad G1--G4 (Firmas Independientes)}
\textbf{\textcolor{warnred}{Sin firma correspondiente no se energiza.}} Una ejecución de LTspice exitosa no representa aprobación de G2, G3 o G4.

\par\vspace{0.1cm}\noindent
\renewcommand{\arraystretch}{1.38}
{\small
\begin{tabularx}{\linewidth}{|c|>{\raggedright\arraybackslash}X|>{\raggedright\arraybackslash}p{4.8cm}|}
\hline
\rowcolor{ucbgray}
\textbf{Puerta} & \textbf{Requisitos de Verificación en Frío} & \textbf{Estado, Firma y Fecha} \\ \hline
\textbf{G1 / E1} & Ficha real NMOS y pinout/pestaña; $100\,\Omega$ potencia; source/tierras de fuente-generador-osciloscopio; control verificado sin bus; limitación de fuente; $R_G/R_{PD}$; OFF previo; instrumentos/corriente apropiados. & Aut.\ / Obs.\ / Bloq.: \rule{1.5cm}{0.4pt}\newline Firma: \rule{2.7cm}{0.4pt}\newline Fecha: \rule{2.7cm}{0.4pt} \\ \hline
\textbf{G2 / E2} & Todo G1 aplicable + \textbf{separación A/B y \texttt{PC817/7805} comprobados} + BJT exactos/pinout, corriente de base y colector, disipaciones, driver y niveles, startup OFF, pérdida de señal OFF, alimentación/desacoplo y prueba previa escalonada sin potencia (Fases 0--D); netlist y log; cero sobreimpulso sobre $V_{GS,\max}$. & Aut.\ / Obs.\ / Bloq.: \rule{1.5cm}{0.4pt}\newline Firma: \rule{2.7cm}{0.4pt}\newline Fecha: \rule{2.7cm}{0.4pt} \\ \hline
\textbf{G3 / E3} & Todo G2 + motor/bobina real, corriente y energía inductiva, diodo \texttt{DF} polaridad/rating, picos de tensión, limitación y forma de medir corriente, temperatura y desconexión. & Aut.\ / Obs.\ / Bloq.: \rule{1.5cm}{0.4pt}\newline Firma: \rule{2.7cm}{0.4pt}\newline Fecha: \rule{2.7cm}{0.4pt} \\ \hline
\textbf{G4 / E4} & Todo G3 pertinente + cuatro MOSFET y drivers físicos, pull-ups P/high-side, referencia segura de $V_{GS}$ high-side, layout/potencia, enclavamiento y dead-time verificados, tabla de estados, no inversión insegura, motor/bus/fusibles, apagado y fallo seguro. & Aut.\ / Obs.\ / Bloq.: \rule{1.5cm}{0.4pt}\newline Firma: \rule{2.7cm}{0.4pt}\newline Fecha: \rule{2.7cm}{0.4pt} \\ \hline
\end{tabularx}
}

\section*{9. Entrega Final, Criterios y Responsabilidades}
Entregar \textbf{un solo informe incremental del Hito 3}, no cuatro informes inconexos, incluyendo: BOM final con componentes realmente utilizados y costos \textit{solo comprobados}, cuatro esquemas/versiones (incluyendo inventario diferencial de migración v1.3$\rightarrow$v1.5), cálculo de diseño, archivos y capturas LTspice con modelos y logs, autorizaciones G1--G4, métodos de medida y capturas reales, registro de problemas y modificaciones, tabla de control de puente H, pérdidas y estimación térmica/Python, valoración crítica de discrepancias y respuestas individuales cortas. Si falta hardware o autorización, presentar bloqueo trazable (\texttt{componente/modelo/medida faltante}, responsable, propuesta de sustitución y evidencia existente); nunca inventar curvas, corrientes, temperaturas ni certificaciones.

\noindent\textbf{Preguntas individuales mínimas:} (1) ¿Qué diferencia existe entre $V_{GS(\text{th})}$ y gate de baja pérdida? (2) ¿Qué significa $R_{DS(\text{on})}$ condicionado? (3) ¿Por qué un seguidor de emisor no convierte $5\,\text{V}$ en $12\,\text{V}$? (4) ¿Qué corriente circula al apagar $L$? (5) ¿Cuál es el riesgo de gate PMOS a $5\,\text{V}$ con source $12\,\text{V}$? (6) ¿Cómo se evita solapamiento en un brazo H? (7) ¿Qué resultado es calculado, simulado o medido?

\section*{10. Fuentes y Estado de Validación}
\begin{itemize}[noitemsep, topsep=2pt]
    \item \textbf{MOSFETs Infineon:} IRLZ44N (\url{https://www.infineon.com/part/IRLZ44N}), IRFZ44N (\url{https://www.infineon.com/part/IRFZ44N}), IRF4905 (\url{https://www.infineon.com/assets/row/public/documents/24/49/infineon-irf4905-datasheet-en.pdf}).
    \item \textbf{Regulador, Optoacoplador y BJTs:} STMicroelectronics L78 series (\url{https://www.st.com/resource/en/datasheet/l78.pdf}); onsemi P2N2222A (\url{https://www.onsemi.com/pdf/datasheet/p2n2222a-d.pdf}) y PN2907A (\url{https://www.onsemi.com/download/data-sheet/pdf/pn2907a-d.pdf}); \texttt{PC817} requiere ficha exacta del fabricante y rango de CTR; folleto SKMI SK3000 aportado (MAIN $50\,\Omega$, sin offset ajustable).
    \item \textbf{Sin verificar todavía:} Inventario/costos exactos BJTs y PMOS, polarización real de \texttt{Q0/Q2} con todas las variantes, respuesta CTR en peor caso, modelo de motor/corriente inductiva, diodo E3, adaptación de high-side, dead-time cuantificado y resultados físicos medidos.
\end{itemize}

\end{document}
